
\subsection{Grupės $\pinp$ ir $\pinm$}

Algebros $\mathup{Cl^\pm(V)}$ poaibis sudarytas iš elemtų $\vb v_1 \vb v_2
\dots \vb v_N$, kur visi $\vb v_i$ tenkina $(\vb v_i,\vb v_i) = 1$, kartu su
algebros daugyba sudaro atitinkamą grupę $\pinpm$. Grupės uždarumo sąvybė akivaizdžiai tenkinama, 
vienetinis elementas yra tas pats kaip ir algebros vienetinis elementas $\mathbb 1$,
o kaip gaunamas atvikštinis elementas siriasi grupeių $\pinp$ ir $\pinm$.

Grupėje $\pinp$ elemento $v$ atvirkštinis elementas yra tranponuotas $v^t$
\begin{multline}
	\vb v_1 \vb v_2 \dots \vb v_N (\vb v_1 \vb v_2 \dots \vb v_N )^t
	=
	\vb v_1 \vb v_2 \dots \vb v_N \vb v_N \vb v_{N-1} \dots \vb v_1
	=
	(\vb v_N,\vb v_N)
	(\vb v_{N-1},\vb v_{N-1})
	\dots
	(\vb v_1,\vb v_1)
	\mathbb 1
	=
	\mathbb 1
\end{multline}
Grupėje $\pinm$ atvirkštinis elementas yra $\alpha(v^t)$
\begin{multline}
	\vb v_1 \vb v_2 \dots \vb v_N \alpha\qty(\vb v_1 \vb v_2 \dots \vb v_N )^t
	=
	\vb v_1 \vb v_2 \dots \vb v_N \alpha(\vb v_N) \alpha(\vb v_{N-1}) \dots \alpha(\vb v_1)
	\\=
	\vb v_1 \vb v_2 \dots \vb v_N (-\vb v_N) (-\vb v_{N-1}) \dots (-\vb v_1)
	=
	(\vb v_N,\vb v_N)
	(\vb v_{N-1},\vb v_{N-1})
	\dots
	(\vb v_1,\vb v_1)
	\mathbb 1
	=
	\mathbb 1
\end{multline}

Grupė $\spin$ galima gauti dviem būdais: kaip $\pinp$ pogrupį sudarytą
iš lyginių vektorių sandaugų arba kaip $\pinm$ pogrupį gautą taip pat.
Abi konstrukcijos yra viena kitai izomofiškos, nes minusas iš sandaugos
$\pinm$ pasirodys lyginį kartą sakičių ir rezultatas bus toks pat kaip
$\pinp$ pogrupyje.

\subsection{$\spin$ ir $\pinpm$ veikimas ant $\set R^3$}

Visos aptartos grupės veikia vektorius iš $\set R^3$, juos pirma patalpinant į
$\mathup{Cl^\pm(V)}$ ir tada su algebros sandauga
\begin{equation}
	\rho(v).x = \begin{cases}
		vxv^t		& v \in \pinp \\
		\alpha(v)xv^t	& v \in \pinm 
	\end{cases}
\end{equation}



